\documentclass[12pt,a4paper]{article}

\usepackage{indentfirst}
\usepackage[utf8]{inputenc}
\usepackage[T1]{fontenc}
\usepackage[brazil]{babel}
\usepackage{amsmath}
\usepackage{geometry}
\usepackage{setspace}
\usepackage{listings}
\usepackage{xcolor}
\usepackage{setspace}
\usepackage{graphicx}
\usepackage{float}
\usepackage{caption}
\usepackage[colorlinks=true,
            linkcolor=blue,
            citecolor=blue,
            urlcolor=blue]{hyperref}

\definecolor{keywordcolor}{rgb}{0,0,0.6}
\definecolor{commentcolor}{rgb}{0,0.5,0}
\definecolor{stringcolor}{rgb}{0.6,0,0}
\definecolor{backgroundcolor}{rgb}{0.95,0.95,0.95}

\lstset{
    language=Matlab,
    backgroundcolor=\color{backgroundcolor},
    basicstyle=\ttfamily\small,
    keywordstyle=\color{keywordcolor}\bfseries,
    commentstyle=\color{commentcolor}\itshape,
    stringstyle=\color{stringcolor},
    numbers=left,
    numberstyle=\tiny\color{gray},
    stepnumber=1,
    numbersep=8pt,
    frame=none,
    rulecolor=\color{black},
    breaklines=true,
    showstringspaces=false
}

\setstretch{1.2}
\setlength{\parindent}{0pt}
\setlength{\parskip}{0.2\baselineskip}
\newcommand{\MATLAB}{\textsc{Matlab}}

\geometry{margin=2.5cm}
\onehalfspacing

\lstset{
  basicstyle=\ttfamily\small,
  frame=single,
  breaklines=true
}

\setlength{\parindent}{2em}

\begin{document}

\begin{titlepage}
    \begin{center}
		\Huge{Universidade de São Paulo}\\
		\large{Escola de Engenharia de São Carlos}\\ 
        %\vspace{.5cm}
        \large{SEM0530 - Problemas de Engenharia Mecatrônica II}\\ 
        \vspace{5cm}
		\Large{Modelagem da rigidez de uma estrutura unidimensional}\\ 
		\vspace{0.3cm}
    
    \end{center}
\vspace{5cm}
	\begin{flushleft}
		
		Aluno: Enzo Larger Manfio  Nº USP: 16863802\\
        Docente: Prof. Dr. Kayc Wayhs Lopes\\
	
    \end{flushleft}
	
	\vfill
	\begin{center}
		\vspace{5cm}
			 São Carlos, Abril de 2026
	\end{center}
\end{titlepage}

\newpage
\tableofcontents


\newpage
\section{Introdução}

Na análise de estruturas frequentemente é feita a modelagem destas como conjuntos de elementos discretos. Deste modo, fenômenos mecânicos podem ser estudados como o agregado de fenômenos individuais mais simples, com o uso de métodos computacionais otimizados. Nesta prática, será evidenciado um exemplo deste modo de análise a partir do estudo da deformação da estrutura unidimensional representada na figura \ref{estrutura}, uma vez sujeita a forças axiais.
\begin{figure}[h]
    \centering
    \includegraphics[width=0.5\textwidth]{estrutura.pdf}
    \caption{Estrutura a ser analisada}
    \label{estrutura}
\end{figure}
\par Mais especificamente, a estrutura será modelada como um número arbitrário de elementos discretos, de modo a encontrar uma expressão generalizada para a Matriz de Rigidez que
relaciona os deslocamentos de cada elemento com as forças aplicadas à estrutura.
Por fim, será desenvolvido um programa na linguagem de programação Matlab que, a
partir de certos parâmetros da estrutura, determine a sua Matriz de Rigidez e calcule o
deslocamento dos elementos uma vez sujeitos a determinadas forças.
\vspace{0.5\baselineskip}

\section{Metodologia}

\subsection{Modelagem matemática}
\vspace{0.5\baselineskip}

\par A Estrutura será modelada como sendo um conjunto de n molas em série, como representado na figura \ref{molas}.
\begin{figure}[h]
    \centering
    \includegraphics[width=0.5\textwidth]{molas.pdf}
    \caption{Estrutura representada como conjunto de n molas em série}
    \label{molas}
\end{figure}
Para cada mola, a sua constante elástica $k_i$ é dada por \eqref{equação k}.
\begin{equation}
\label{equação k}
    k_n = k_{min} + \Delta{}e^{-bi}
\end{equation}
\par O i-ésimo ponto analisado está sujeito à uma força externa correspondente, $F_i$, e as forças elásticas das molas que ligam o ponto i-1 à i e i à i+1,
Com isso, considerando o ponto em equilíbrio e o sentido das forças e deslocamentos na figura \ref{estrutura}, temos:
\begin{math}
\begin{center}
    \sum{F_x} = 0 \\ \nonumber
    F_i + F_{Ki+1} - F_{Ki} = 0 \\ \nonumber
    F_i + k_{i+1}(u_{i+1} - u_k) - k_i(u_i - u_{k-1}) = 0 \\\nonumber
\end{center}
\end{math}
\begin{equation}
    \label{Equação Força}
    F_i = -k_iu_{i-1} + (k_i + k_{i +1})u_i - k_{i+1}u_{i+1}
\end{equation}
 \par Em que $u_i$ e $k_i$ são deslocamento e constante elástica da mola correspondente ao i-ésimo elemento. Com isso, é possível definir um sistema matricial representando a relação entre os
 vetores $\textbf{u} = \{u_1, u_2,  ... u_n\}^T$ e $\textbf{F} = \{F_1, F_2, ... F_n\}^T$:

 \begin{equation}
    \label{Sistema}
    \textbf{Mu} = \textbf{F}
 \end{equation}
Em que
\begin{equation}
    \label{MatrizRigidez}
    \textbf{M} =
    \begin{bmatrix}
    k_1 + k_2 & -k_2 & 0 & \cdots & 0 \\
    -k_2 & k_2 + k_3 & - k_3 &\cdots & 0 \\
    0 & -k_3& k_3 + k_4 &\cdots & 0 \\
    \vdots & \vdots & \vdots & \ddots & \vdots \\
    0 & 0 & 0 & \cdots & k_n
    \end{bmatrix}
\end{equation}

\vspace{0.5\baselineskip}

\noindent constitui a Matriz de Rigidez do sistema, a qual, por ser simétrica, é necessariamente inversível, de modo que o sistema \eqref{Sistema} possui solução única.

\vspace{0.5\baselineskip}

\subsection{Implementação no \MATLAB}

\par Em um arquivo próprio, foi definida a função VetorKi(), que recebe como argumento parâmetros os parâmetros gerais do sistema de molas ($k_{min}$ e b) e o número de molas, retornando o vetor das constantes elásticas das molas, calculadas com \eqref{equação k}.

\begin{lstlisting}[language=Matlab, 
caption = {Função para criação do vetor de constantes elásticas.},
label = {lst:funcao_ki}]
function k_i = VetorKi(n, k_min, delta_k, b)
    k_i = zeros(n, 1);
    for i = 1:n
        k_i(i) = k_min + delta_k*exp(-b*i);
    end
end
\end{lstlisting}

\par Em um outro arquivo, foi definida a função MatrizRigidez() para, a partir do vetor de constantes de elasticidade, construir a matriz de rigidez correspondente ao sistema, utilizando \eqref{Equação Força} para cada linha da matriz.
\begin{lstlisting}[language=Matlab, 
caption = {Função para criação da matriz de rigidez},
label = {lst:funcao_matriz}]
function Matriz = MatrizRigidez(k)
    n = length(k);
    Matriz = zeros(n, n);
    Matriz(1, 1) = k(1);
    Matriz(n, n) = k(n);
    if n>1
        Matriz(1, 1) = k(1) + k(2);
        Matriz(1, 2) = -k(2);
        Matriz(n, n-1) = -k(n);
    end
    if n > 2
        for i = 2:n-1
            Matriz(i, i-1) = -k(i);
            Matriz(i, i) = k(i) + k(i+1) ;
            Matriz(i, i + 1) = -k(i+1);
        end
    end
end
\end{lstlisting}
\par Por fim, utilizou-se a função linsolve() oferecida pelo Matlab para resolver o sistema \eqref{Sistema} e foram plotados os gráficos de $k_i$ e $u_i$ em função de i.
\section{Resultados}

\noindent Foi considerada a estrutura sob duas configurações:
\begin{itemize}
\item 30 molas, uma força de 50 N é aplicada na extremidade livre.
\item 10 molas, uma força de -50 N é aplicada no meio da estrutura (i = 5) e uma de 100 N é aplicada na extremidade livre.
\end{itemize}
Para ambas, considerou-se $k_{min} = 30$ kN/m, b = 0,15 e $\Delta{k} = (70 + 0.2\cdot{02})$ kN/m.
\subsection{Configuração 1}

\begin{figure}[H]
    \centering
    \includegraphics[width=0.7\textwidth]{Figure_1.pdf}
    \caption{Rigidez e deslocamento de cada elemento para a configuração 1}
    \label{grafico 1}
\end{figure}
\par Como o deslocamento é inversamente proporcional à rigidez sob força constante (Caso da Configuração 1, já que a força aplicada na extremidade é transmitida ao longo da estrutura), conforme a rigidez de cada elemento decresce exponencialmente o deslocamento de cada elemento aumenta (já que considerou-se o sentido da tração como positivo).
\subsection{Configuração 2}

\begin{figure}[H]
    \centering
    \includegraphics[width=0.7\textwidth]{Figure_2.pdf}
    \caption{Rigidez e deslocamento de cada elemento para a configuração 2}
    \label{grafico 2}
\end{figure}

\par É possível observar que o gráfico dos deslocamentos até i = 5 é parecido com o da configuração 1. Isso ocorre pois, somando as forças de tração de 100 N e de compressão de 50 N, a secção correspondente à primeira metade do corpo está efetivamente tracionada por 50 N, assim como a estrutura na configuração 1. A partir de i = 5, a estrutura não está mais sujeita à essa força de compressão, resultando em um aumento no crescimento do deslocamento.

\section{Conclusão}

\par Neste trabalho, foi demonstrada a modelagem de estruturas mecânicas a partir de sua discretização com o estudo de uma estrutura unidimensional como um conjunto de molas em série. Uma vez considerado o equilíbrio de cada extremidade livre das molas, foi possível encontrar uma expressão geral para a Matriz de Rigidez que relaciona o deslocamento destas e as forças aplicadas.
\par Por fim, foram desenvolvidos scripts em Matlab que calculam a matriz de rigidez e resolvem o sistema linear correspondente para determinar os deslocamentos de pontos da estrutura uma vez aplicadas forças sobre estes, sendo testadas duas configurações para a estrutura, com os resultados obtidos para a deformação desta sendo fisicamente consistentes em ambas.
\end{document}